% chapters/05/section-03.tex
\section{Gerak Berbasis Selang Waktu}
\label{sec:gerak-selang-waktu}

Inilah bagian inti bab ini: menghubungkan kecepatan yang telah disusun pada
bagian sebelumnya dengan waktu yang telah diukur pada Bab~IV.

\subsection{Perpindahan Selama Satu Bingkai}

Perpindahan objek selama satu bingkai dihitung dari kecepatan dan selang
waktu. Selang waktu yang dipakai adalah yang benar-benar
berlangsung, yaitu nilai yang dikembalikan \code{clock.tick()}.

\begin{kodekeluaran}
velocity = direction.normalize() * speed
position = position + velocity * delta_seconds
\end{kodekeluaran}

Baris kedua itulah yang membuat kecepatan objek tidak bergantung pada laju
bingkai. Kecepatan dinyatakan dalam piksel per detik, selang waktu dinyatakan
dalam detik, sehingga hasil kalinya adalah piksel. Apabila selang waktunya
dua kali lebih panjang, perpindahannya dua kali lebih jauh, dan jarak yang
ditempuh dalam satu detik tetap sama.

Perhatikan urutan perhitungannya. Kecepatan dihitung lebih dahulu, baru
posisi diperbarui. Urutan itu sama dengan kaidah Euler yang telah dijelaskan
pada Bagian~2.4: memperbarui besaran turunan lebih dahulu, baru besaran yang
diturunkan darinya.

\subsection{Membuktikan Kecepatan Tetap melalui Pengukuran}

Pernyataan bahwa kecepatan objek tetap sama pada setiap laju bingkai sebaiknya
dibuktikan, bukan hanya dipercaya. Program pada bab ini menyediakan alat
pembuktiannya: setiap bingkai, program mencatat perubahan posisi yang
sesungguhnya, lalu membagi jumlah jarak dengan lama objek benar-benar
bergerak.

Cara itu penting. Apabila jaraknya dihitung dari kecepatan yang ditetapkan,
hasilnya pasti sama dengan kecepatan itu sendiri, dan pengukurannya tidak
membuktikan apa pun. Yang diukur di sini adalah perubahan posisi yang nyata.

Tabel~\ref{tab:pengukuran-kecepatan} memuat hasil pengukuran pada satu
komputer contoh, dengan kecepatan yang ditetapkan 200 piksel per detik dan
objek bergerak ke kanan.

\begin{table}[htbp]
\centering
\caption{Hasil pengukuran kecepatan pada tiga laju bingkai}
\label{tab:pengukuran-kecepatan}
{\setlength{\tabcolsep}{4pt}\footnotesize%
\begin{tabularx}{\textwidth}{@{}C{1.8cm} C{1.5cm} C{2.6cm} C{2.4cm} Y@{}}
\toprule
\headrow \textbf{Batas laju} & \textbf{Bingkai} & \textbf{Waktu berjalan} & \textbf{Jarak tempuh} & \textbf{Kecepatan terukur} \\
\midrule
30 FPS & 18 & 0,61 detik & 115,4 piksel & 200,0 piksel/detik \\
60 FPS & 36 & 0,60 detik & 116,8 piksel & 200,0 piksel/detik \\
120 FPS & 72 & 0,62 detik & 122,4 piksel & 200,0 piksel/detik \\
\bottomrule
\end{tabularx}}
\tabsource{Pengukuran penulis pada kecepatan 200 piksel per detik, arah lurus,
tanpa jendela. Angka pada komputer Anda akan berbeda; yang perlu Anda
perhatikan adalah kolom terakhir.}
\end{table}

Kolom terakhir menunjukkan hal yang penting: kecepatan terukur sama persis
pada ketiga laju bingkai. Pada 30 bingkai per detik objek berpindah lebih jauh
setiap bingkainya, sedangkan pada 120 bingkai per detik perpindahannya lebih
pendek, tetapi jarak yang ditempuh dalam satu detik sama. Inilah bukti bahwa
gerak tidak lagi bergantung pada laju bingkai.

\begin{catatabox}
Perhatikan kolom jarak tempuh: nilainya berbeda-beda (115,4 sampai 122,4
piksel) walaupun kecepatannya sama. Selisih itu bukan kesalahan.
\code{clock.tick()} hanya menjamin laju tidak \textbf{melampaui} batas, bukan
bahwa setiap bingkai berlangsung tepat sepanjang jatahnya; waktu berjalan pada
tabel di atas memang berbeda-beda. Tambahan pula, selang waktu bingkai
pertama bernilai nol sehingga bingkai itu belum menghasilkan perpindahan.
Yang harus sama adalah \textbf{kecepatan}, bukan jarak tempuh selama selang
waktu yang diukur secara kasar.
\end{catatabox}

\subsection{Mengapa Cara Ini Perlu Diuji Sendiri}

Mahasiswa yang telah membaca bab ini sampai di sini sering menganggap
pembuktian tersebut berlebihan, karena ``jelas kecepatannya sama''. Anggapan
itu berbahaya:

\begin{itemize}
  \item Apabila \code{delta_seconds} lupa dipakai sebagai pengali, program
        tetap berjalan dan tidak menimbulkan pesan kesalahan apa pun. Yang
        terlihat hanyalah gerak yang berbeda-beda menurut komputer. Tanpa
        pengukuran, kekeliruan itu dapat bertahan sampai akhir semester.
  \item Apabila \code{delta_seconds} dikalikan dua kali, misalnya karena
        kecepatan sudah dinyatakan dalam satuan per bingkai, objek tetap
        bergerak dan hanya terasa terlalu cepat. Gejala itu mudah disalahkan
        kepada nilai kecepatan yang dipilih.
  \item Apabila penormalan vektor arah terlupa, gerak diagonal tetap tampak
        wajar pada permainan sederhana. Kekeliruannya baru terasa ketika
        pemain menyadari bahwa berlari diagonal lebih cepat daripada berlari
        lurus, dan pemain berpengalaman akan segera menyadarinya.
\end{itemize}

Karena itu, biasakan menjalankan program gerak pada dua laju bingkai yang
berbeda setiap kali selesai menulisnya, misalnya 30 dan 60. Apabila kecepatan
objek terasa berbeda, ada yang keliru pada perhitungan perpindahannya. Cara
yang sama akan Anda pakai terus pada bab-bab berikutnya \cite{fiedler-timestep}.
